\section{Expected Outputs and Contributions}

The expected outcome is that OPSD will improve VQA performance relative to conventional 
SFT in cases in which the privileged cropped region provides stronger visual evidence than the full-image.
This expected outcome is based on the findings of \textcite{yuan2026vision} that 
multimodal models could perform better when conditioned on evidence-centered crops versus full images.

However, the size of the improvement is expected to vary based on the model size and architecture.
Each model may display different levels of regional-to-global perception gaps. As a result, different
amounts of useful knowledge may be available for self-distillation. This would allow the study to provide 
evidence on how dependent the effectiveness of OPSD is on VLM architecture and scale.

The key contribution of this project would be the controlled comparison between OPSD and SFT for
general-domain Visual Genome VQA across multiple VLM architectures and scales. Potential challenges 
include the additional computational cost needed for each teacher and student forward pass as well as the 
possibility that some image crops may provide little to no advantage over the original full size image. 